\documentclass[12pt,a4paper]{article}

% Define \noopsort for bibliography sorting (used for names like "van de Kuilen")
%\providecommand{\noopsort}[1]{}

% utf8 support
\usepackage[utf8]{inputenc}
% PDF font support
\usepackage[T1]{fontenc}
% lmodern fonts
\usepackage{lmodern}
\usepackage{eurosym}
\usepackage{textcomp}
% double spacing
\usepackage[nodisplayskipstretch]{setspace}
% Ragged right with indentation
\usepackage[document]{ragged2e}
\setlength{\RaggedRightParindent}{\parindent}
% Indent the first paragraph
\usepackage{indentfirst}

% Captions for figures
\usepackage[justification=centering]{caption}
\usepackage{subcaption}

\usepackage{geometry}
\usepackage{xcolor}

% For quoting
\usepackage[style = american]{csquotes}
%To make references into links
\usepackage[hidelinks]{hyperref}
% For equation stuff
\usepackage{mathtools}
% For xfrac
\usepackage{xfrac}
% For math symbols
\usepackage{amsmath}
\usepackage{amssymb}
\usepackage[mathcal]{euscript}
% For Bold math Symbols
\usepackage{bm}
% For nice looking fractions
\usepackage{nicefrac}
% For Graphics
\usepackage{graphicx}
% To float pictures appropriately
\usepackage{float}
% Read about this, tested it, and was sold.
\usepackage[stretch=15]{microtype}

% Adjust titles
% and section rules
\usepackage{titling}
\usepackage[raggedright]{titlesec}
\titleformat*{\subsection}{\large\itshape}

% Include PDF pages directly
\usepackage{pdfpages}
% For landscape pages
\usepackage{pdflscape}

% Stuff for Tables
\usepackage{tabularx}
\newcolumntype{Y}{>{\centering\arraybackslash}X} % A centered column type that sucks up excess space
% For \toprule, \midrule, and \bottomrule
\usepackage{booktabs}
% Have a value over multiple rows in a table
\usepackage{multirow}

% Squish things onto the page
\usepackage{adjustbox}

% PythonTeX for embedding Python code
\usepackage{pythontex}

% bibliography stuff
\usepackage[american]{babel}
\usepackage[
	backend=biber,
	bibstyle=IEEEtran,
	uniquename=false,
	doi=false,
	uniquelist=false,
	isbn=false,
	url=false,
	style=authoryear-comp,
	maxcitenames=1,
	maxcitenames=4,
	maxbibnames=999,
	useprefix=true
]{biblatex}
% No Oxford Commas
\DefineBibliographyExtras{english}{\let\finalandcomma=\empty}
% Don't put "in:" in anything
\renewbibmacro{in:}{}
% The path to the bibliography file
\IfFileExists{/home/bam/Dropbox/Papers/library.bib}{
	\addbibresource{/home/bam/Dropbox/Papers/library.bib}
}{}
\addbibresource{./library.bib}
% Add some space between entries
\setlength\bibitemsep{1.5\itemsep}
% Define \noopsort for Dutch name sorting (e.g., van de Kuilen sorted under K)
\providecommand{\noopsort}[1]{}

% Convenience Commands
\newcommand{\overbar}[1]{\mkern 1.5mu\overline{\mkern-1.5mu#1\mkern-1.5mu}\mkern 1.5mu} % Command to make a shorter overline/longer bar

% Make a break a column title into 2 lines
\newcommand{\cnline}[2][c]{%
	\setlength{\extrarowheight}{0em}
	\begin{tabular}[#1]{@{}c@{}}#2\end{tabular}
	\setlength{\extrarowheight}{5em}
}

%% Personal Commands
\newcommand\email[1]{\href{mailto:#1}{#1}}

% Sincgle space the sections
\let\oldsection\section
\renewcommand{\section}[2][]{\singlespacing\oldsection#1{#2}\doublespacing}
\let\newsection\section

\let\oldsubsection\subsection
\renewcommand{\subsection}[1]{\singlespacing\oldsubsection{#1}\doublespacing}
\let\newsubsection\subsection


% Some commands to help with appendicies
\let\pbib\printbibliography
\renewcommand{\printbibliography}{

	\let\section\oldsection
	\let\subsection\oldsubsection

	\singlespacing
	\pbib
	\doublespacing

	\let\section\newsection
	\let\subsection\newsubsection
}

\newcommand{\appendixsetup}{
	\appendix
	\titleformat{\section}[hang]
		{\normalfont\normalsize\bfseries}
		{Appendix \thesection}
		{0pt}
		{: \normalsize\bfseries\filcenter}
	\renewcommand{\thetable}{\Alph{section}.\arabic{table}}
	\renewcommand{\thefigure}{\Alph{section}.\arabic{figure}}
	\renewcommand{\theequation}{\Alph{section}.\arabic{equation}}
}

\newcommand{\appendixref}{
	\titleformat{\section}[hang]
		{\normalfont\normalsize\bfseries}
		{Appendix \thesection}
		{0pt}
		{\normalsize\bfseries\filcenter}

	\singlespacing
	\printbibliography
	\doublespacing

	\titleformat{\section}[hang]
		{\normalfont\normalsize\bfseries}
		{Appendix \thesection}
		{0pt}
		{: \normalsize\bfseries\filcenter}
}

\newcommand{\appsection}[1]{
	\renewcommand{\thetable}{\Alph{section}.\arabic{table}}
	\renewcommand{\thefigure}{\Alph{section}.\arabic{figure}}
	\renewcommand{\theequation}{\Alph{section}.\arabic{equation}}
	\renewcommand{\theHequation}{\Alph{section}.\arabic{equation}}
	\setcounter{table}{0}
	\setcounter{figure}{0}
	\setcounter{equation}{0}
	\section{#1}
}

%% Commands for figures and captions

\newcommand{\addfig}[2][]{
  \IfFileExists{#2}{
    \includegraphics[#1]{#2}
  }{
    \textbf{FIGURE: \enquote{\detokenize{#2}} NOT FOUND! Go make it in the code!}
  }
}

% Let the captions be actually centered!
\newcommand{\mycaption}[1]{
  \setlength{\RaggedRightParindent}{0pt}
  \caption{#1}
  \setlength{\RaggedRightParindent}{\parindent}
}

\newcommand{\mycaptionof}[2]{
  \setlength{\RaggedRightParindent}{0pt}
  \captionof{#1}{#2}
  \setlength{\RaggedRightParindent}{\parindent}
}



\hyphenpenalty=10000

\title{Principled Hyperprior Selection for Bayesian Hierarchical Models of Risk Preferences}
\author{
	Brian Albert Monroe\\
    \small{Center for the Economic Analysis of Risk, }\\
    \small{Maurice R. Greenberg School of Risk Science,}\\
    \small{J. Mack Robinson College of Business,}\\
    \small{Georgia State University}\\
		\small{bmonroe3@gsu.edu}}


\newcommand\blfootnote[1]{
    \begingroup
    \renewcommand\thefootnote{}\footnote{#1}
    \addtocounter{footnote}{-1}
    \endgroup
}

% Math shortcuts
\newcommand{\E}[1]{\mathbb{E}\left[#1\right]}
\newcommand{\Var}[1]{\text{Var}\left[#1\right]}
\newcommand{\Cov}[2]{\text{Cov}\left[#1, #2\right]}

\begin{document}

% Load generated data from NPZ files
\begin{pycode}
import numpy as np
from pathlib import Path

# Find data directory relative to doc/
doc_dir = Path('.').resolve()
data_dir = doc_dir.parent / 'data' / 'output'

def _load_npz(name):
    f = data_dir / name
    return dict(np.load(f)) if f.exists() else {}

def fmt(x, decimals=2):
    """Format number, handling numpy types."""
    return f"{float(x):.{decimals}f}"

# Load all NPZ files
D = _load_npz('paper_data.npz')
Q = _load_npz('paper_quantities.npz')
_prior_params = _load_npz('prior_params.npz')
_ce_comp = _load_npz('posterior_ce_comparison.npz')
_ce_overlay = _load_npz('predictive_ce_overlay.npz')

# Total posterior samples for inference
total_samples = f"{int(Q.get('n_samples_post_burnin', 0)) * int(Q.get('n_chains', 0)):,}"

# Subject CE identifiers
subject_most_affected = int(_ce_comp['most_affected_sid']) if _ce_comp else 0
subject_least_affected = int(_ce_comp['least_affected_sid']) if _ce_comp else 0
subject_max_ce = subject_most_affected
subject_min_ce = subject_least_affected

# CE overlay quantities
hprior_pred_rs = fmt(float(_ce_overlay.get('prior_pct_above_ev', 0)), 1)
prior_ce_below_15 = fmt(float(_ce_overlay.get('prior_pct_below_15', 0)), 1)
prior_ce_above_40 = fmt(float(_ce_overlay.get('prior_pct_above_40', 0)), 1)
info_ce_below_15 = fmt(float(_ce_overlay.get('info_pct_below_15', 0)), 1)
info_ce_above_40 = fmt(float(_ce_overlay.get('info_pct_above_40', 0)), 1)
uninf_ce_below_15 = fmt(float(_ce_overlay.get('uninf_pct_below_15', 0)), 1)
uninf_ce_above_40 = fmt(float(_ce_overlay.get('uninf_pct_above_40', 0)), 1)
\end{pycode}


\maketitle
\pagenumbering{gobble}
\singlespacing{}

\begin{abstract}
  \begin{justify}

  Bayesian hierarchical models (BHMs) are increasingly used to estimate
  heterogeneous risk preferences, yet researchers face the difficult task of
  specifying hyperprior distributions over population parameters. I propose a
  principled approach: specify priors for a single random individual---a
  cognitively tractable task---then derive hyperpriors such that the
  unconditional BHM prior matches this individual-level specification. I
  demonstrate the approach using rank-dependent utility estimation with data
  from \textcite{Harrison2016}, showing that the methodology produces
  coherent hyperpriors from interpretable inputs and well-identified
  posteriors. A comparison with a diffuse hyperprior illustrates how
  hyperprior choice can have substantive consequences for individual-level
  inference.

  \end{justify}
\end{abstract}
\vspace{\fill}
\small
\textbf{Keywords:} Bayesian hierarchical models, risk preferences, hyperprior selection, rank-dependent utility

\textbf{JEL Classification:} C11, C25, D81

\normalsize
\clearpage
\pagenumbering{arabic}
\setcounter{page}{1}
\doublespacing

\section{Introduction}

Bayesian hierarchical models (BHMs) are becoming a standard tool for estimating
heterogeneous preferences in experimental economics \parencite{Gao2022}. These
models jointly estimate individual-level parameters and the hyper-parameters
that govern their prior distributions. With the assumption that individuals are
exchangeable, BHMs have the advantage of using the distribution of individual
estimands to inform the prior distribution for the estimands. There is no free
lunch, however: instead of needing a specified prior distribution for a given
individual, one must now specify a \textit{hyperprior} over the
\textit{hyper-parameters} of the individual level estimands.

It is common for models to assign a prior distribution with 2 hyper-parameters
per individual-level parameter; for example, a normal distribution with a mean
and variance. For BHMs, each hyper-parameter needs its own prior distribution,
which is again often defined with 2 degrees of freedom; for example, another
normal distribution for the hyper-mean parameter and an inverse-gamma
distribution for the hyper-variance parameter. Thus, a model with $p$
parameters typically results in a hyperprior that requires the specification of
$4p$ hyperprior parameters.

Common models of risk preferences modeled with Expected Utility Theory (EUT)
can have 2-3 parameters, models assuming Rank Dependent Utility (RDU) can have
3-5 parameters, and models assuming Cumulative Prospect Theory (CPT) can have
up to 10.\footnote{
  These parameter count ranges assume all models share a noise parameter, and
  use utility and probability weighting functions. Two common utility functions
  are the CRRA and Expo-Power \parencite{Saha1993} functions, with 1 and 2
  parameters each, giving EUT 2-3 parameters total. Two popular probability
  weighting functions are the power function \parencite{Quiggin1982}, and the
  two parameter function by \textcite{Prelec1998}, giving RDU models 3-5
  parameters. CPT has a utility and probability weighting functions separately
  defined over losses and gains, leading to a base of 5 parameters for a simple
  parameterization with the CRRA and Power functions and one noise parameter,
  up to 10 parameters with Expo-Power and Prelec functions and separate noise
  parameters for losses and gains.
} These parameters interact in complex ways, making the choice of picking 4
hyperprior parameters per model parameter difficult.

In this paper, I propose a principled way to select hyperprior parameters for
BHMs of risk preferences by first specifying a prior predictive distribution
for a single agent, and then, noting the assumed exchangeability of agents,
deriving the hyperprior that produces the same distribution. Since it is often
simpler and more natural to think of risk preferences for individual agents, I
believe this process results in proper hyperprior distributions that can help
researchers pick informed priors, and help critics of research using BHMs focus
on priors of convenience that ignore substantive knowledge about the target
domain.

I begin by outlining the process for selecting a hyperprior for an arbitrary
model. I then demonstrate the approach using rank-dependent utility (RDU)
estimation with lottery choice data from \textcite{Harrison2016}. I show that
the method produces well-identified posteriors and illustrate how hyperprior
choice can have substantive consequences for individual-level inference.

\section{The Prior Selection Approach}

I begin with an arbitrary, single parameter model:
\begin{equation}
  \label{eq:simple_model}
  y_i = \Phi(\theta_i, x),
\end{equation}
\noindent where $y_i$ is some modeled quantity conditional on other observable
data, $x$, and $\theta_i \in \mathbb{R}$ is a parameter of interest for a single
subject. A Bayesian model seeks to find the posterior distribution of
$\theta_i$ conditional on the observed data:
\begin{align}
  \label{eq:bayes_post}
  \begin{split}
    p(\theta_i | y,x) &= \frac{p(y | \theta_i, x) p(\theta_i)}{\displaystyle\int p(y | \theta_i, x) p(\theta_i) \, d\theta_i} \\
    {} &\propto p(y | \theta_i, x) p(\theta_i) \\
  \end{split}
\end{align}
\noindent where $p(\theta_i | y,x)$ is the posterior of the parameter of
interest, $p(y | \theta_i, x)$ the likelihood of the modeled quantity
conditional on $\theta_i$ and the observed data $x$, and $p(\theta_i)$ is the
prior distribution of $\theta_i$.

While the prior distribution $p(\theta_i)$ can be any proper probability
distribution,\footnote{
  And in some cases, an improper distribution.
} it is useful for now to think of $\theta_i \in \mathbb{R}$ and normally
distributed:
\begin{equation}
\label{eq:simple_prior}
  \theta_i \sim \text{Normal}(\mu, \sigma^2)
\end{equation}
\noindent where $\mu$ and $\sigma^2$ are prior mean and variance parameters
signifying the researcher's prior belief about the distribution of possible
$\theta_i$ parameters that describe an individual.

If we have data from multiple individuals, it is reasonable to think that
different posterior distributions of $\theta_i$ across subjects may influence
our beliefs about the population they were drawn from. This influence is
directly modeled with a Bayesian Hierarchical Model (BHM), whereby the
prior parameters are themselves said to be distributed according to some
\textit{hyperprior} distribution.
\begin{align}
  \label{eq:simple_bhm_mean}
  \mu &\sim \text{Normal}(\mu_0, \tau_0^2) \\
  \label{eq:simple_bhm_sd}
  \sigma^2 &\sim \text{Inverse-Gamma}(\alpha_0, \beta_0)
\end{align}
\noindent The use of Normal and Inverse-Gamma distributions is common for
real-valued hierarchical models because it lends itself to the conjugate
relationship of normal distributions.\footnote{
  For now, I assume the model takes a single parameter, but it should be clear
  that a set of uncorrelated parameters $\theta^j_i$, with $j \in [1, p]$ can
  be specified in the same way. Additionally, while I initially assume
  independence across all parameters $\theta^j_i$, the case with full
  covariance structure is discussed in Appendix \ref{app:full_cov}. For now, I
  drop the $j$ superscript for convenience.
}

The simple BHM model specified in equations
(\ref{eq:simple_prior}-\ref{eq:simple_bhm_sd}) requires the estimation of 4
hyper-parameters for parameter $\theta_i$. Directly specifying these values can be
challenging because their joint implications are often difficult to assess.
What does it mean to set $\alpha_{0} = 3$ versus $\alpha_{0} = 10$? How should
$\tau_{0}^2$ relate to $\beta_{0}$? These questions often lack intuitive
answers, even if $\theta_i$ is well understood.

\subsection{Exchangeability and a prior over an individual}

The prior selection approach begins with a cognitively tractable question:
\emph{What might the value $\theta_i$ of a single random individual look like?}
This requires revisiting the problem of selecting a prior in
(\ref{eq:simple_prior}), and the direct specification of $\mu$ and $\sigma^2$
prior parameters. This lower-level specification of prior parameters is often
conceptually easier than the direct specification of hyperprior parameters. Let
us designate this \enquote{direct prior} as:
\begin{equation}
  \label{eq:direct_prior}
  p(\theta_i | \mu^*, \sigma^{2*}) = \text{Normal}(\mu^*, \sigma^{2*}) \\
\end{equation}
\noindent Note here that we use the superscript ${}^*$ to designate our direct prior
parameters. When omitted, the parameter is modeled rather than specified.

While the question posed asks about an individual, not a population, its answer
implicitly integrates over population heterogeneity introduced in
(\ref{eq:simple_bhm_mean})-(\ref{eq:simple_bhm_sd}):
\begin{equation}
  \label{eq:prior_restriction}
    p(\theta_i |\  \mu^*, \sigma^{2*}) = \int p(\theta_i |\  \mu, \sigma^2) \, p(\mu, \sigma^2 |\  \boldsymbol{\psi} ) \, d\mu \, d\sigma^2
\end{equation}
\noindent where $\boldsymbol{\psi}$ denotes the vector of hyperprior
parameters, whatever distributional family is chosen for $p(\mu, \sigma^2)$.
This restriction, that a prior specified over a given individual implicitly
integrates over all possible heterogeneity that could have produced that prior,
is the key component to the proposed hyperprior selection process. It amounts to
a consistency requirement between beliefs about the hyperprior parameters
$\boldsymbol{\psi}$ and beliefs about what a random individual would look like
directly, $\mu^*$ and $\sigma^{2*}$.

The solution to picking the hyperprior parameters is reconciling these two
sets of parameters. If a direct prior has two degrees of freedom, like a normal
distribution, then any hyperprior must have:
\begin{align}
\label{eq:prior_decomposed}
\begin{split}
  \E{\theta_i} &= \E{\E{\theta_i | \mu}} \\
  \Var{\theta_i} &= \E{\sigma^2} + \Var{\mu}
\end{split}
\end{align}
\noindent Where the first part of (\ref{eq:prior_decomposed}) is due to the law
of iterated expectations, and the second part is due to the law of total
variance.\footnote{
  In general, if the direct prior had $n$ degrees of freedom, then
  (\ref{eq:prior_decomposed}) would require the expansion of $n$ equations
  using the law of total cumulance. For cumulant $\kappa_n$:
  \begin{equation*}
    \kappa_n(Y) = \E{\kappa_n(Y|X)} + g_n(\kappa_1, \ldots, \kappa_{n-1})
  \end{equation*}
  where $g_n$ involves lower-order cumulants. For $n=1$ (mean), $\kappa_1(Y) =
  \E{\kappa_1(Y|X)}$, which is the law of iterated expectations. For $n=2$
  (variance), $\kappa_2(Y) = \E{\kappa_2(Y|X)} + \kappa_2(\kappa_1(Y|X))$,
  which is the law of total variance. An example using the 3 parameter
  skew-Normal distribution is given in Appendix \ref{app:skew_normal}.
}
Importantly, (\ref{eq:prior_restriction}) and (\ref{eq:prior_decomposed}) hold
for \textit{any} hyperprior distribution with defined first and second
moments---they are consequences of the laws of iterated expectations and total
variance, not properties of a particular distributional family. The method
developed below therefore applies generally; only the final mapping from moment
targets to specific hyperparameters depends on the chosen family.

\subsection{Mean and Variance Decomposition}

The first part of (\ref{eq:prior_decomposed}) is relatively simple for a normal
distribution given its symmetry:
\begin{align}
\begin{split}
\mu^* &= \E{\theta_i} \\
  {} &= \E{\E{\theta_i | \mu}} \\
  {} &= \E{\mu} \\
  {} &= \mu_0
\end{split}
\end{align}
\noindent Thus, the mean of the $\mu$ \textit{hyper-parameter}, must be equal
to the mean of the direct prior.

The decomposition of the variance parameter, however, requires some additional
re-parameterization. The law of total variance gives:
\begin{equation}
    \label{eq:total_var}
    \sigma^{2*} = \Var{\theta_i} = \E{\sigma^2} + \Var{\mu}
\end{equation}
\noindent This decomposition allows us to think of the direct prior variance as
the sum of two components: uncertainty about the population mean, $\Var{\mu}$,
and expected population heterogeneity, $\E{\sigma^2}$. We can re-parameterize
this relationship by assigning a fraction of the direct prior variance
$\sigma^{2*}$ to each component:
\begin{align}
  \label{eq:rho_var_mu}
  \rho \sigma^{2*} &= \Var{\mu} \\
  \label{eq:rho_E_sigma}
  (1 - \rho) \sigma^{2*} &= \E{\sigma^2}
\end{align}
\noindent where $\rho \in (0,1)$ is a variance partition ratio: the fraction of
direct prior variance attributable to uncertainty about where the population is
centered. Together,
(\ref{eq:prior_restriction})--(\ref{eq:rho_E_sigma}) constitute the general
framework. They hold for any hyperprior family and reduce the problem to
specifying $\rho$ and matching the moments of the chosen family to these
targets.

For the Normal-Inverse-Gamma specification in
(\ref{eq:simple_bhm_mean})--(\ref{eq:simple_bhm_sd}), $\Var{\mu} = \tau_0^2$
and $\E{\sigma^2} = \beta_0/(\alpha_0 - 1)$, so the moment targets become:
\begin{align}
  \rho \sigma^{2*} &= \tau_{0}^2 \\
  (1 - \rho) \sigma^{2*} &= \frac{\beta_0}{\left(\alpha_0 - 1\right)}
\end{align}
\noindent By specifying $\rho$ and a known direct prior variance $\sigma^{2*}$,
we fully define $\tau_0^2$.

The expected population variance $\E{\sigma_j^2} = \beta_{0,j}/(\alpha_{0,j} -
1)$, however, involves two unknowns. We therefore either need to further
re-parameterize, or specify one directly. With the $\rho$ parameter in place,
and interpretable in part as managing prior beliefs about the heterogeneity
of the population, specifying $\alpha_0$ directly is a reasonable
course of action. The $\alpha_0$ parameter is the shape of the Inverse-Gamma
distribution that determines the variance hyper-parameter. In a sense,
$\alpha_0$ can be thought of as a pseudo-count of the number of observations
that informs the prior.

The $\rho$ parameter can be interpreted as the central
tendency of our beliefs about the heterogeneity of the sample, with low values
of $\rho$ indicating a greater belief in a high degree of heterogeneity, and
high values of $\rho$ indicating the opposite. The choice of the $\alpha_0$
parameter can likewise be interpreted as expressing our \textit{confidence} in
that degree of heterogeneity. Low values of $\alpha$ would result in a wider
distribution of the $\sigma^2$ hyper-parameter with heavy right tails, while
greater values of $\alpha$ would rein in those tails, making the distribution
of $\sigma^2$ more heavily concentrated around its mean.

\textit{Any} choice of $\rho \in [0, 1]$ and $\alpha > 0$ will result in the
unconditional hyperprior being equivalent to the direct prior. Additionally,
the specification of $\rho$ and $\alpha$ is \textit{unit-agnostic}. Neither
require intuition about the possible values of the mean or variance
hyperparameter. That unit-specific intuition about possible values is baked
into the specification of the direct prior and the imposed restriction. This is
very different from having to specify the hyperprior directly, which would
require unit-specific intuition about $\tau_0$, the variance of the mean
hyperparameter, and $\beta$, the scale of the variance hyperparameter.

The choice the researcher needs to make, after having specified the direct
prior, is therefore an expression of their belief about the degree of
heterogeneity in the sample. While it could be the case that we have different
beliefs about this heterogeneity composition per-parameter, it may also be a
reasonable default to set the same $\rho$ and $\alpha$ hyper-prior parameters
for all parameters.

With the researcher selection of a direct prior, $\mu^*$ and $\sigma^{2*}$, and
the selection of the variance partition ratio and shape of population variance,
$\rho$ and $\alpha_0$, we have 4 specified parameters, matching the initial 4
degrees of freedom needed to select hyperprior parameters. With these
specifications, the hyper-parameters specified in (\ref{eq:simple_bhm_mean})
and (\ref{eq:simple_bhm_sd}) are fully defined:
\begin{align}
    \mu_{0} &= \mu^* \label{eq:hp_mu}\\
    \tau_{0}^2 &= \rho \sigma^{2*} \label{eq:hp_tau}\\
    \alpha_0 &= \alpha_0 \quad \text{(researcher specified)} \label{eq:hp_alpha}\\
    \beta_{0} &= (1 - \rho) \sigma^{2*} (\alpha_0 - 1) \label{eq:hp_beta}
\end{align}

It is worth noting that there is no \enquote{free lunch} here; with 4
hyperprior parameters, there are still 4 decisions needed to satisfy the 4
degrees of freedom of the hyperprior. This approach is most advantageous when
it is easier to specify the prior for a given individual $\theta_i$ through a
simpler distribution and separately consider how sample heterogeneity would
present in that direct sample, than to specify the 4 hyperprior parameters
directly.

\section{Defining a Bayesian Hierarchical Model of Risk Preferences}

To demonstrate this method I define a BHM where all subjects have RDU risk
preferences using the CRRA utility function and the 2-parameter
\textcite{Prelec1998} probability weighting function (PWF).
\begin{align}
  \label{eq:rdu}
  U  &= \sum_i d_i u(x_i)\\
  \label{eq:crra}
  u(x) &= \frac{x^{1-r}}{1 - r}
\end{align}
\noindent where $d_i$ is the decision weight associated with outcome $x_i$. For
a lottery with $n$ outcomes ranked from worst ($x_1$) to best ($x_n$), let
$p_i$ denote the probability of outcome $x_i$. The 2-parameter
\textcite{Prelec1998} probability weighting function transforms probabilities
into decisions weights as follows:
\begin{align}
  c_i &= \sum_{j=1}^{i} p_j \\
  \label{eq:prelec}
  w_i &= \exp\left(-\eta \left(-\ln c_i\right)^\phi\right)\\
  \label{eq:decision_weight}
  d_i &= w_i - w_{i - 1}
\end{align}
with $w_0 = 0$ and $\phi, \eta > 0$. When $\phi = \eta = 1$, the probability
weighting function is linear, resulting in $d_i = p_i$, the special case of
EUT\@.

The likelihood that a lottery \textit{A} is chosen over \textit{B} follows the
logistic choice function:
\begin{align}
\label{eq:likelihood}
\begin{split}
  L(A) &= \frac{1}{1 + \exp\left(-\Delta\right)}, \\
  \Delta &= \frac{U(A) - U(B)}{\lambda \left(u(\bar{x}) - u(\underbar{x})\right)}
\end{split}
\end{align}
\noindent where $\Delta$ is a non-linear latent index of utility differences,
$\lambda > 0$ is a \enquote{Fechnerian} noise scaling parameter, and
$u(\bar{x}) - u(\underbar{x})$ is the difference between the greatest and least
outcome with positive probability across \textit{A} and \textit{B}. Dividing by
the difference of maximum and minimum utilities normalizes the difference in
utilities to be between $[-1, 1]$ and applies the Contextual Utility model of
\textcite{Wilcox2011}.

This basic model has been estimated numerous times in varied applications. It
was used by \textcite{Harrison2021} to assess the changes in risk preferences
of during the COVID-19 pandemic, it was used to recover subjective beliefs in
\textcite{Harrison2019a}, and many more applications. \textcite{Gao2022} make
use of this model in a BHM to demonstrate how individual preferences can be
informed by the assumption of exchangeability among subjects.

As a model for an individual agent, the vector $\theta_i = [r_i, \phi_i,
\eta_i, \lambda_i]$ must be estimated. In a Bayesian framework, this estimation
takes the form of finding the posterior distribution of $\theta$, conditional
on observed choice data:
\begin{equation}
  p(\theta_i | D) \propto L(D | \theta_i) p(\theta_i)
\end{equation}
\noindent where (\ref{eq:likelihood}) specifies how the likelihood $L(D |
\theta_i)$ is calculated, and the job of the researcher is to specify the prior
distribution $p(\theta_i)$.

\subsection{Prior Specification of a Single RDU Agent}

\begin{pycode}
# Build structured dicts from already-loaded NPZ data
dprior = {
    "r": {"mu": float(D.get("dprior_r_mu", 0)),
          "sigma2": float(D.get("dprior_r_sigma2", 0))},
    "phi_raw": {"mu": float(D.get("dprior_phi_raw_mu", 0)),
                "sigma2": float(D.get("dprior_phi_raw_sigma2", 0))},
    "eta_raw": {"mu": float(D.get("dprior_eta_raw_mu", 0)),
                "sigma2": float(D.get("dprior_eta_raw_sigma2", 0))},
    "lambda_raw": {"mu": float(D.get("dprior_lambda_raw_mu", 0)),
                   "sigma2": float(D.get("dprior_lambda_raw_sigma2", 0))}
}

ps_rho = float(D.get("ps_rho", 0))
ps_alpha_0 = float(D.get("ps_alpha_0", 0))

hyperpriors = {
    "r": {
        "mu_0": float(_prior_params.get("r_mean", 0)),
        "tau_0": float(np.exp(_prior_params.get("ln_r_sdev", 0))),
        "alpha_0": float(_prior_params.get("alpha_0", 0)),
        "beta_0": float(_prior_params.get("beta_0_r", 0)),
    },
    "phi": {
        "mu_0": float(_prior_params.get("phi_mean", 0)),
        "tau_0": float(np.exp(_prior_params.get("ln_phi_sdev", 0))),
        "alpha_0": float(_prior_params.get("alpha_0", 0)),
        "beta_0": float(_prior_params.get("beta_0_phi", 0)),
    },
    "eta": {
        "mu_0": float(_prior_params.get("eta_mean", 0)),
        "tau_0": float(np.exp(_prior_params.get("ln_eta_sdev", 0))),
        "alpha_0": float(_prior_params.get("alpha_0", 0)),
        "beta_0": float(_prior_params.get("beta_0_eta", 0)),
    },
    "lambda": {
        "mu_0": float(_prior_params.get("lambda_mean", 0)),
        "tau_0": float(np.exp(_prior_params.get("ln_lambda_sdev", 0))),
        "alpha_0": float(_prior_params.get("alpha_0", 0)),
        "beta_0": float(_prior_params.get("beta_0_lambda", 0)),
    },
}
\end{pycode}

Following the example above, I define the $r_i$ parameter, and the natural log
of the remaining parameters, to have a Normal distribution with specific values
for these \enquote{direct prior} parameters:
\begin{equation}
\label{eq:d_prior}
\begin{matrix}
  r \sim N\left(\mu^{r} = \py{dprior["r"]["mu"]},\right.& \left.\sigma^{2, r} = \py{dprior["r"]["sigma2"]}\right) \\
  \ln\phi \sim N\left(\mu^{\phi} = \py{dprior["phi_raw"]["mu"]},\right.& \left.\sigma^{2, \phi} = \py{dprior["phi_raw"]["sigma2"]}\right) \\
  \ln\eta \sim N\left(\mu^{\eta} = \py{dprior["eta_raw"]["mu"]},\right.& \left.\sigma^{2, \eta} = \py{dprior["eta_raw"]["sigma2"]}\right) \\
  \ln\lambda \sim N\left(\mu^{\lambda} = \py{dprior["lambda_raw"]["mu"]},\right.& \left.\sigma^{2, \lambda} = \py{dprior["lambda_raw"]["sigma2"]}\right) \\
\end{matrix}
\end{equation}
\noindent These prior choices represent my a-priori belief about what a
reasonable subject may look like. The median subject is modestly risk averse,
and conforms to EUT\@. However, for the purposes of the current discussion, any
values could have been chosen.

Using the method described above, in addition to the two direct prior
parameters per parameter in \ref{eq:d_prior}, I need to characterize my
beliefs about the heterogeneity in the sample by specifying the $\rho$ and
$\alpha_0$ hyperpriors for each parameter. I believe there to be a fair degree
of heterogeneity in the population, but I am also not dogmatically certain
about the location of the mean hyperparameter. I therefore pick a $\rho$ value
closer to 0 than to 1, and an $\alpha_0$ parameter indicating that my
confidence comes from the equivalent of having seen roughly \py{ps_alpha_0}
subjects. I do not have great reason to think heterogeneity differs
systematically across parameters, so I set the same values for all parameters:
\begin{align}
  \rho^j = \py{ps_rho} \quad &\forall j \in [r, \phi, \eta, \lambda] \\
  \alpha_{0}^j = \py{ps_alpha_0} \quad &\forall j \in [r, \phi, \eta, \lambda]
\end{align}
\noindent These researcher decisions result in the fully specified hyperprior:
\begin{align}
  \label{eq:hyperprior-values}
\begin{matrix}
\mu^r \sim N\left(\mu_0^r = \py{f"{hyperpriors['r']['mu_0']:.2f}"}, \tau_0^r = \py{f"{hyperpriors['r']['tau_0']:.2f}"}\right), & \sigma^{2, r} \sim IG\left(\alpha_0^r = \py{ps_alpha_0}, \beta_0^r = \py{f"{hyperpriors['r']['beta_0']:.2f}"}\right) \\[4pt]
\mu^\phi \sim N\left(\mu_0^\phi = \py{f"{hyperpriors['phi']['mu_0']:.2f}"}, \tau_0^\phi = \py{f"{hyperpriors['phi']['tau_0']:.2f}"}\right), & \sigma^{2, \phi} \sim IG\left(\alpha_0^\phi = \py{ps_alpha_0}, \beta_0^\phi = \py{f"{hyperpriors['phi']['beta_0']:.2f}"}\right) \\[4pt]
\mu^\eta \sim N\left(\mu_0^\eta = \py{f"{hyperpriors['eta']['mu_0']:.2f}"}, \tau_0^\eta = \py{f"{hyperpriors['eta']['tau_0']:.2f}"}\right), & \sigma^{2, \eta} \sim IG\left(\alpha_0^\eta = \py{ps_alpha_0}, \beta_0^\eta = \py{f"{hyperpriors['eta']['beta_0']:.2f}"}\right) \\[4pt]
\mu^\lambda \sim N\left(\mu_0^\lambda = \py{f"{hyperpriors['lambda']['mu_0']:.2f}"}, \tau_0^\lambda = \py{f"{hyperpriors['lambda']['tau_0']:.2f}"}\right), & \sigma^{2, \lambda} \sim IG\left(\alpha_0^\lambda = \py{ps_alpha_0}, \beta_0^\lambda = \py{f"{hyperpriors['lambda']['beta_0']:.2f}"}\right) \\[4pt]
\end{matrix}
\end{align}


%\begin{figure}
%  \centering
%  \mycaption{Draws from the Direct Prior}
%  \label{fig:direct_prior}
%  \includegraphics[width=\textwidth]{../plots/priors/direct_prior.pdf}
%\end{figure}
%
%\begin{figure}
%  \centering
%  \mycaption{Draws from the Derived hyperprior}
%  \label{fig:}
%  \includegraphics[width=\textwidth]{../plots/priors/unconditional_prior.pdf}
%\end{figure}
%
%\begin{figure}
%  \centering
%  \mycaption{Draws from the Derived hyperprior}
%  \label{fig:}
%  \includegraphics[width=\textwidth]{../plots/priors/prior_comparison.pdf}
%\end{figure}



\section{Results}


I use the risk preference task from \textcite{Harrison2016}, which elicited
lottery choices from \py{D['n_subjects']} subjects in a laboratory experiment
in 2015 at Georgia State University. Each subject made
\py{D['n_choices_per_subject']} binary choices between three-outcome lotteries
with prize amounts of \$10, \$30, and \$50. I estimate the model using a Markov
Chain Monte Carlo (MCMC) with \py{Q.get('n_chains', 'N/A')} parallel
chains, ultimately obtaining \py{total_samples} samples from the posterior for
inference.\footnote{
  The sampler uses a block-wise Metropolis-Hastings in Gibbs sampler with
  adaptive proposals informed by cross-chain covariance estimation. The maximum
  $\hat{R}$ statistic is \py{fmt(Q.get('rhat_max', 0), 2)} across all
  \py{Q.get('n_total_params', 'N/A')} parameters (\py{Q.get('n_subjects',
  'N/A')} subjects $\times$ 4 parameters + 8 hyperparameters), suggesting
  excellent convergence. Further details of the sampling procedure along with
  trace plots of the hyperparameters, and individual-level parameters for the
  subsequently discussed subjects \py{subject_min_ce} and \py{subject_max_ce},
  are provided in Appendix \ref{app:mcmc}.
}

I estimate the model using the informative hyperprior derived in
(\ref{eq:hyperprior-values}). To illustrate the consequences of hyperprior
choice, I also estimate the model using the diffuse hyperprior from
\textcite[279-281]{Gao2022}:
\begin{align}
  \label{eq:ght_prior}
\begin{matrix}
\mu^r \sim N\left(\mu_0^r             = 0, \tau_0^r = 100\right), & \sigma^{2, r} \sim IG\left(\alpha_0^r = 0.001, \beta_0^r = 0.001\right) \\[4pt]
\mu^\phi \sim N\left(\mu_0^\phi       = 0, \tau_0^\phi =100\right), & \sigma^{2, \phi} \sim IG\left(\alpha_0^\phi = 0.001, \beta_0^\phi = 0.001\right) \\[4pt]
\mu^\eta \sim N\left(\mu_0^\eta       = 0, \tau_0^\eta =100\right), & \sigma^{2, \eta} \sim IG\left(\alpha_0^\eta = 0.001, \beta_0^\eta = 0.001\right) \\[4pt]
\mu^\lambda \sim N\left(\mu_0^\lambda = 0, \tau_0^\lambda =100\right), & \sigma^{2, \lambda} \sim IG\left(\alpha_0^\lambda = 0.001, \beta_0^\lambda = 0.001\right) \\[4pt]
\end{matrix}
\end{align}

I present results in two ways. First, I examine the posterior distributions of
the hyperparameters under each specification, showing how the informative
hyperprior contributes to posterior identification. Second, I examine the
posterior predictive distributions of a certainty equivalent (CE) defined over
an equi-probable lottery, illustrating how hyperprior choice affects
individual-level inference.

\subsection{Hyperparameter Posteriors}

\begin{figure}[hp]
  \centering
  \mycaption{Hyperparameter posteriors}
  \label{fig:hyperparam_post}
  \includegraphics[width=\textwidth]{../plots/paper/hyper_comparison.pdf}
\end{figure}
\begin{figure}[hp]
  \centering
  \mycaption{Hyper-posterior distribution}
  \label{fig:hyperparam_pred}
  \includegraphics[width=\textwidth,height=0.85\textheight,keepaspectratio]{../plots/paper/predictive_params_overlay.pdf}
\end{figure}

Figure \ref{fig:hyperparam_post} displays the posterior densities of the
hyperparameters under the informative prior (blue), the diffuse prior
(orange), and the informative hyperprior itself (grey).

The informative hyperprior contributes additional information beyond the
likelihood, concentrating the posterior distributions relative to the diffuse
case for all hyperparameters except $\mu_\lambda$. The posteriors under
the informative hyperprior are also generally closer to the hyperprior mode,
reflecting the expected influence of the prior on posterior location, with
the exception of $\mu_\phi$. These parameters cannot be interpreted fully
independently, however; the posterior of the $\sigma_\phi$ term is much
closer to the hyperprior mode, which may compensate.

Figure \ref{fig:hyperparam_pred} shows the implied unconditional density for
$\theta_i = [r_i, \phi_i, \eta_i, \lambda_i]$ under three cases: the
informative hyperprior (prior predictive), the informative posterior, and the
diffuse posterior. The figure tells a similar story to Figure
\ref{fig:hyperparam_post} but in more interpretable terms: the informative
posterior remains close to the prior predictive distribution, while the
diffuse posterior diverges more substantially. Notably, although the
informative posterior of the mean and variance hyperparameters for $\eta$ are
clearly different from the hyperprior in Figure \ref{fig:hyperparam_post},
the informative posterior and prior predictive distributions for $\eta$ in
Figure \ref{fig:hyperparam_pred} are almost identical.

This can be reconciled through the variance decomposition used to derive the
hyperprior. If the two unconditional distributions are similar but the
variance on the mean hyperparameter (characterized by $\tau_0$) is lower for
the posterior than the hyperprior, this implies an effective $\rho$ closer to
0 for the posterior---interpretable as greater population heterogeneity. This
is exactly what Figure \ref{fig:hyperparam_post} shows: the posterior mean of
$\ln \sigma_\eta$ under the informative prior is greater than the
corresponding hyperprior, reflecting more inter-subject heterogeneity.


\subsection{Posterior Predictive Distribution of Certainty Equivalent}

The Certainty Equivalent (CE) of a lottery is the amount of money for certain that
provides an equivalent utility as a given lottery. The CE can be used as a
normalized utility function \parencite{King1983}, but is also very useful in
terms of assessing risk preference models with many interrelated preference
parameters \parencite{Harrison2022}. The CE for a lottery with outcomes $x \in
[10, 30, 50]$ and $p(x) = 1/3$ for all $x$,\footnote{
  These outcomes represent the 3 outcomes in the dataset of \textcite{Harrison2016}.
} is defined as:
\begin{align}
\label{eq:CE}
  \begin{split}
  \frac{\text{CE}^{1 - r}}{1 - r} &= \sum_x d_i(x) \frac{x^{1 - r}}{1 - r} \\
    \text{CE} &= \left(\sum_x d_i(x) x^{1 - r}\right)^\frac{1}{1 - r}
  \end{split}
\end{align}

\begin{figure}[h]
  \centering
  \mycaption{Hyper-posterior predictive CE distributions}
  \label{fig:hyper_post_ce}
  \includegraphics[width=\textwidth]{../plots/paper/predictive_ce_overlay.pdf}
\end{figure}

The CE in (\ref{eq:CE}) can be calculated for any vector $\theta_i = [r_i,
\phi_i, \eta_i, \lambda_i]$. Using the sampled draws of $\theta_i$ from the
posterior distribution results in a posterior predictive distribution of CEs. I
focus on two results with respect to the CE. Figure \ref{fig:hyper_post_ce}
shows the hyperprior predictive distribution of the CE using the informed
hyperprior, and the hyper-posterior predictive of the CE for both informative
and uninformative priors. Figures \ref{fig:post_ce_min} and
\ref{fig:post_ce_max} examine two individual subjects predictive posterior
distributions. Figure \ref{fig:post_ce_min} shows the posterior predictive
distribution of the subject who has the smallest difference in the median CE
between the informative and uninformative models, while Figure
\ref{fig:post_ce_max} shows the subject with the largest difference in median
CE.

Figure \ref{fig:hyper_post_ce} provides a broader picture of how hyperprior
choice affects posterior inferences about risk preferences. The informative
hyperprior implies \py{hprior_pred_rs}\% of prior predictive CE values
greater than the expected value of the lottery, reflecting a not insignificant
prior belief that a random subject might be risk seeking. Of greater interest,
however, are the posterior predictive distributions.

The diffuse posterior predictive distribution, shown in orange, displays
a wide distribution of possible CEs, with substantial density at extreme values
of risk aversion. The mode and \py{uninf_ce_below_15}\% of the density of the
diffuse posterior is less than \$15, indicating extreme risk aversion.
Likewise, \py{uninf_ce_above_40}\% of the density is above \$40, indicating
extreme risk \textit{seeking}. In contrast, only \py{info_ce_below_15}\% of the
CE density is below \$15 for the informative posterior and
\py{info_ce_above_40}\% above \$40. The informative hyperprior concentrates
the posterior predictive density away from these extreme values.

\begin{figure}[hp]
  \centering
  \mycaption{Posterior predictive CE: Least Affected (Subject \py{subject_min_ce})}
  \label{fig:post_ce_min}
  \includegraphics[width=\textwidth]{../plots/paper/subject_ce_least_affected.pdf}
\end{figure}
\begin{figure}[hp]
  \centering
  \mycaption{Posterior predictive CE: Most Affected (Subject \py{subject_max_ce})}
  \label{fig:post_ce_max}
  \includegraphics[width=\textwidth]{../plots/paper/subject_ce_most_affected.pdf}
\end{figure}

Figure \ref{fig:hyper_post_ce}, however, only shows the hyper-posterior
predictive distribution. With a BHM, I have posterior distributions of
$\theta_i$ for each subject and can therefore examine each of the 111 subjects'
predictive posterior distribution directly. I focus for the moment on two
subjects. Subject \py{subject_min_ce}, displayed in Figure
\ref{fig:post_ce_min} has the smallest change in median CE between the
informative and uninformative models, and subject \py{subject_max_ce},
displayed in Figure \ref{fig:post_ce_max}, has the largest. Both subjects'
posterior distributions show interesting, and potentially counter-intuitive
results.

Starting with Figure \ref{fig:post_ce_min}, there is effectively no
difference in the medians of the two posterior distributions, and both
posteriors are substantially more peaked than the prior distribution. This
suggests that the data for this subject were very informative with regard to
the RDU model, perhaps that the subject made very
consistent choices. When the data are very informative, it makes sense that the
prior distribution would have less weight in the posterior distribution.
Perhaps counter-intuitively, the posterior of the informative model is
\textit{more peaked} than the uninformative model, despite the prior
distribution being flatter than the posterior of the uninformative model. At
first glance this may be odd, but what is not shown in this figure is prior
predictive distribution of the \textit{uninformative} model, which is
substantially flatter than the informative prior in the figure.

In Figure \ref{fig:post_ce_max}, I show the subject with the greatest
difference in median predictive posterior distributions between the two
specifications. The posterior densities are much flatter than for the
previous subject, indicating that the data for this subject were much less
informative. The effect of the hyperprior is therefore much more pronounced.
The median of the posterior predictive distribution under the informative
hyperprior indicates near risk neutrality, with roughly half of the density
on CE values that indicate risk seeking. Under the diffuse hyperprior, the
posterior predictive distribution differs substantially, illustrating how
hyperprior choice can materially affect inference when data are less
informative.

What both figures show is that the inclusion of an informative hyperprior
matters most when the data are not particularly informative about parameters of
interest, which is very intuitive, but also that an informative hyperprior
interacts at the individual subject level only indirectly.


\section{Discussion}


While it is often the case that picking hyperprior parameters can be difficult
and unintuitive, re-parameterization of the problem, and recognition of the
exchangeability principle, can aid a researcher who wants to do more than
specify \enquote{uninformative} hyperpriors. The path chosen in this paper, in
the context of a BHM of risk preferences, is to recognize that researchers
often have strong beliefs about what a typical subject's preferences might look
like, and use these individual-level beliefs to inform population-level
beliefs. These beliefs are often informed by decades of published empirical
results, and expert intuition about what constitutes reasonable and probable
preferences. By recognizing that under exchangeability, and without any ex-ante
data to the contrary, the belief about what a typical subject might look like
must apply to \textit{all} subjects in a sample \textit{a-priori}, I establish the
primary result of this paper: the unconditional \enquote{direct} prior about an
individual integrates over the uncertainty about the population. This limiting
condition on hyperprior distributions has not been commonly utilized in BHMs.

This restricting condition does not result in a \enquote{free lunch}, however.
The same degrees of freedom must ultimately be satisfied. To this end, I
propose a variance decomposition that has the useful property of being
unit-agnostic, and is interpretable in terms of the researcher's belief about
the degree of heterogeneity in their sample. Taken together, the specification
of a direct prior about what a typical subject ought to look like, and the
expression of belief about the degree of heterogeneity in the sample, the
proposed method resolves all degrees of freedom and does so in a way that
allows for useful hyperpriors to be specified by rigorous researchers, and
scrutinized by those who may worry that prior distributions are being abused to
produce arbitrary posterior inferences.

The application illustrates that hyperprior choice has consequences for
inference, particularly for subjects whose data are less informative---and
that the proposed method provides a coherent, interpretable way to make that
choice. What I have not discussed is how to pick the direct prior of a
typical individual, or how to compare two models with the same data and
likelihood functions but different hyperpriors.

Starting with the latter point, one of the most common ways Bayesian models are
compared is by use of Bayes' factor:
\begin{equation}
  \text{BF} = \frac{m_\text{info}}{m_\text{noinfo}} = \frac{\displaystyle\int p(D | \theta) \pi_\text{info}(\theta) \, d\theta}{\displaystyle\int p(D | \theta) \pi_\text{noinfo}(\theta) \, d\theta}
\end{equation}
\noindent But in this case, the only difference between the two models is the
prior distribution, $\pi$. The Bayes' factor then reduces entirely to a
comparison of how well each prior predicted the data before even seeing the
data. Totally uninformative hyperpriors would lose to even somewhat accurate
prior distributions, even if the posterior inferences are nearly identical, as
they are for subject \py{subject_min_ce} in Figure \ref{fig:post_ce_min}. A
model where the prior and posterior match exactly will always win in the Bayes'
factor horse race, creating a perverse incentive to evaluate the model, then
re-evaluate the model again using the previous posterior as the current prior,
until the point where prior matches posterior. The ability to unscrupulously
\enquote{game} Bayes' factor makes it unappealing as a method to adjudicate
what constitutes a better prior.

Instead the way a critic ought to judge a prior follows the same path that a
rigorous researcher ought to take when trying to define the prior in the first
place: the prior ought to be defended as the true beliefs of the researcher
about the parameters of interest. These beliefs ought to be informed by past
data, expertise, theory, and, sometimes, mathematical convenience. Critics
ought to interrogate chosen priors on the same grounds. The method proposed in
this paper reduces the burden for both parties, the researcher and critic, to
engage in both of their respective responsibilities.

% The References for the main text
\pagebreak
\printbibliography
\pagebreak

% Appendices
\appendixsetup{}

\appsection{Full Covariance Derivation}
\label{app:full_cov}
\titleformat*{\subsection}{\normalsize\itshape}

In the main text, I derived hyperpriors for a model with $p$ independent
parameters, each governed by its own Normal-Inverse-Gamma hyperprior. Here I
extend the derivation to the case where the $p$ parameters have a full
covariance structure, using Normal-Inverse-Wishart (NIW) hyperpriors. This
extension is relevant when the researcher believes that the population-level
parameters may be correlated---for example, that subjects with higher risk
aversion also tend to have lower probability weighting curvature.

\subsection{Setup}

Suppose the direct prior for a single individual is multivariate normal:
\begin{equation}
  \label{eq:mv_direct_prior}
  \boldsymbol{\theta}_i \sim \text{MVN}(\boldsymbol{\mu}^*, \boldsymbol{\Sigma}^*)
\end{equation}
where $\boldsymbol{\theta}_i \in \mathbb{R}^p$ is the parameter vector,
$\boldsymbol{\mu}^* \in \mathbb{R}^p$ is the direct prior mean, and
$\boldsymbol{\Sigma}^*$ is a $p \times p$ positive definite covariance matrix.
Note that when $\boldsymbol{\Sigma}^*$ is diagonal, this reduces to $p$
independent normal distributions, recovering the setting of Section 2.

\subsection{Normal-Inverse-Wishart Hierarchical Structure}

For the hierarchical model with correlated parameters, a natural conjugate
choice is the Normal-Inverse-Wishart (NIW) hyperprior:
\begin{align}
  \label{eq:niw_mu}
  \boldsymbol{\mu} \mid \boldsymbol{\Sigma} &\sim
    \text{MVN}\!\left(\boldsymbol{\mu}_0,\;
    \frac{\boldsymbol{\Sigma}}{\kappa_0}\right) \\[4pt]
  \label{eq:niw_sigma}
  \boldsymbol{\Sigma} &\sim \text{Inverse-Wishart}(\boldsymbol{\Psi}_0, \nu_0)
\end{align}
where $\boldsymbol{\mu}_0 \in \mathbb{R}^p$ is the prior mean for the
population mean, $\kappa_0 > 0$ scales the variance of the population mean
relative to the population covariance, $\boldsymbol{\Psi}_0$ is the $p \times
p$ positive definite scale matrix, and $\nu_0 > p + 1$ is the degrees of
freedom parameter.

The NIW distribution has the useful property that the marginal variance of the
population mean depends on the population covariance through $\kappa_0$. This
coupling captures the intuition that uncertainty about the population mean
should scale with the population spread.

\subsection{Mean Constraint}

By the law of iterated expectations, the unconditional mean of
$\boldsymbol{\theta}_i$ must equal the direct prior mean:
\begin{align}
  \E{\boldsymbol{\theta}_i}
    &= \E{\E{\boldsymbol{\theta}_i \mid \boldsymbol{\mu}}} \notag \\
    &= \E{\boldsymbol{\mu}} \notag \\
    &= \boldsymbol{\mu}_0
\end{align}
Thus, as in the univariate case:
\begin{equation}
  \label{eq:niw_mean_constraint}
  \boldsymbol{\mu}_0 = \boldsymbol{\mu}^*
\end{equation}

\subsection{Covariance Constraint}

The matrix analog of the law of total variance gives:
\begin{equation}
  \label{eq:matrix_total_var}
  \text{Cov}[\boldsymbol{\theta}_i]
    = \E{\text{Cov}[\boldsymbol{\theta}_i \mid \boldsymbol{\mu}, \boldsymbol{\Sigma}]}
    + \text{Cov}[\E{\boldsymbol{\theta}_i \mid \boldsymbol{\mu}}]
\end{equation}

Since $\boldsymbol{\theta}_i \mid \boldsymbol{\mu}, \boldsymbol{\Sigma} \sim
\text{MVN}(\boldsymbol{\mu}, \boldsymbol{\Sigma})$, the two terms are:
\begin{align}
  \E{\text{Cov}[\boldsymbol{\theta}_i \mid \boldsymbol{\mu}, \boldsymbol{\Sigma}]}
    &= \E{\boldsymbol{\Sigma}} \\[4pt]
  \text{Cov}[\E{\boldsymbol{\theta}_i \mid \boldsymbol{\mu}}]
    &= \text{Cov}[\boldsymbol{\mu}]
\end{align}

Substituting the NIW moments, the expected covariance is:
\begin{equation}
  \E{\boldsymbol{\Sigma}}
    = \frac{\boldsymbol{\Psi}_0}{\nu_0 - p - 1}
    \qquad (\text{for } \nu_0 > p + 1)
\end{equation}
and the covariance of the mean is:
\begin{equation}
  \text{Cov}[\boldsymbol{\mu}]
    = \E{\frac{\boldsymbol{\Sigma}}{\kappa_0}}
    = \frac{1}{\kappa_0} \E{\boldsymbol{\Sigma}}
    = \frac{\boldsymbol{\Psi}_0}{\kappa_0(\nu_0 - p - 1)}
\end{equation}

Therefore the unconditional covariance is:
\begin{equation}
  \label{eq:niw_total_cov}
  \boldsymbol{\Sigma}^*
    = \E{\boldsymbol{\Sigma}} + \text{Cov}[\boldsymbol{\mu}]
    = \E{\boldsymbol{\Sigma}} \left(1 + \frac{1}{\kappa_0}\right)
    = \frac{\boldsymbol{\Psi}_0}{\nu_0 - p - 1}
      \cdot \frac{\kappa_0 + 1}{\kappa_0}
\end{equation}

\subsection{Variance Partition with \texorpdfstring{$\rho$}{rho}}

As in the univariate case, I decompose the direct prior covariance into
uncertainty about the population mean and expected population heterogeneity,
using the same variance partition ratio $\rho \in (0, 1)$:
\begin{align}
  \label{eq:niw_cov_mu}
  \text{Cov}[\boldsymbol{\mu}]
    &= \rho \cdot \boldsymbol{\Sigma}^* \\[4pt]
  \label{eq:niw_E_sigma}
  \E{\boldsymbol{\Sigma}}
    &= (1 - \rho) \cdot \boldsymbol{\Sigma}^*
\end{align}

The interpretation is identical to the univariate case: $\rho$ represents the
fraction of direct prior variance attributable to uncertainty about where the
population is centered, while $(1 - \rho)$ represents the fraction due to
population heterogeneity.

From $\text{Cov}[\boldsymbol{\mu}] =
\E{\boldsymbol{\Sigma}}/\kappa_0$, substituting
(\ref{eq:niw_cov_mu})-(\ref{eq:niw_E_sigma}):
\begin{equation}
  \rho \cdot \boldsymbol{\Sigma}^*
    = \frac{(1-\rho) \cdot \boldsymbol{\Sigma}^*}{\kappa_0}
\end{equation}
which gives:
\begin{equation}
  \label{eq:niw_kappa}
  \kappa_0 = \frac{1 - \rho}{\rho}
\end{equation}

From $\E{\boldsymbol{\Sigma}} =
\boldsymbol{\Psi}_0 / (\nu_0 - p - 1)$:
\begin{equation}
  \label{eq:niw_psi}
  \boldsymbol{\Psi}_0
    = (1 - \rho) \cdot \boldsymbol{\Sigma}^* \cdot (\nu_0 - p - 1)
\end{equation}

\subsection{Derived Hyperparameters}

The researcher specifies:
\begin{enumerate}
  \item A direct prior $\boldsymbol{\mu}^*$, $\boldsymbol{\Sigma}^*$ (the
    prior for a single individual).
  \item The variance partition ratio $\rho \in (0, 1)$.
  \item The degrees of freedom $\nu_0 > p + 1$ (analogous to $\alpha_0$ in the
    univariate case, controlling confidence in the prior covariance).
\end{enumerate}

\noindent The NIW hyperprior parameters are then fully determined:
\begin{align}
  \boldsymbol{\mu}_0 &= \boldsymbol{\mu}^*
    \label{eq:niw_hp_mu} \\[3pt]
  \kappa_0 &= \frac{1 - \rho}{\rho}
    \label{eq:niw_hp_kappa} \\[3pt]
  \boldsymbol{\Psi}_0 &= (1-\rho) \cdot \boldsymbol{\Sigma}^* \cdot (\nu_0 - p - 1)
    \label{eq:niw_hp_psi} \\[3pt]
  \nu_0 &= \nu_0 \quad \text{(researcher specified)}
    \label{eq:niw_hp_nu}
\end{align}

\subsection{Verification}

To verify, substitute the derived hyperparameters back into
(\ref{eq:niw_total_cov}):
\begin{align}
  \E{\boldsymbol{\Sigma}} + \text{Cov}[\boldsymbol{\mu}]
  &= \frac{\boldsymbol{\Psi}_0}{\nu_0 - p - 1}
     \left(1 + \frac{1}{\kappa_0}\right) \notag \\[4pt]
  &= \frac{(1-\rho) \boldsymbol{\Sigma}^* (\nu_0 - p - 1)}{\nu_0 - p - 1}
     \left(1 + \frac{\rho}{1 - \rho}\right) \notag \\[4pt]
  &= (1-\rho) \boldsymbol{\Sigma}^*
     \cdot \frac{1}{1 - \rho} \notag \\[4pt]
  &= \boldsymbol{\Sigma}^*
\end{align}
confirming that the unconditional covariance of $\boldsymbol{\theta}_i$
recovers the direct prior covariance $\boldsymbol{\Sigma}^*$.

\subsection{Relationship to the Independent Case}

When $\boldsymbol{\Sigma}^* = \text{diag}(\sigma_1^{2*}, \ldots,
\sigma_p^{2*})$, the off-diagonal elements of $\boldsymbol{\Psi}_0$ are zero,
and the Inverse-Wishart reduces to $p$ independent Inverse-Gamma distributions.
Specifically, the marginal distribution of $\Sigma_{jj}$ under
$\boldsymbol{\Sigma} \sim \text{IW}(\boldsymbol{\Psi}_0, \nu_0)$ with diagonal
$\boldsymbol{\Psi}_0$ is:
\begin{equation}
  \Sigma_{jj} \sim \text{Inverse-Gamma}\!\left(
    \frac{\nu_0 - p + 1}{2},\;
    \frac{[\boldsymbol{\Psi}_0]_{jj}}{2}
  \right)
\end{equation}

Substituting (\ref{eq:niw_hp_psi}):
\begin{equation}
  \Sigma_{jj} \sim \text{Inverse-Gamma}\!\left(
    \frac{\nu_0 - p + 1}{2},\;
    \frac{(1-\rho) \sigma_j^{2*} (\nu_0 - p - 1)}{2}
  \right)
\end{equation}

This recovers the independent-parameter formulas
(\ref{eq:hp_alpha})-(\ref{eq:hp_beta}) with the identification:
\begin{align}
  \alpha_0 &= \frac{\nu_0 - p + 1}{2} \\[3pt]
  \beta_0 &= \frac{(1-\rho) \sigma_j^{2*} (\nu_0 - p - 1)}{2}
    = (1 - \rho) \sigma_j^{2*} (\alpha_0 - 1)
\end{align}
where the last equality uses $\nu_0 - p - 1 = 2(\alpha_0 - 1)$. Similarly,
$\text{Cov}[\boldsymbol{\mu}] = \boldsymbol{\Sigma} / \kappa_0$ reduces to
$\text{Var}[\mu_j] = \sigma_j^2 / \kappa_0$, and integrating over the
Inverse-Gamma yields $\tau_{0,j}^2 = \rho \cdot \sigma_j^{2*}$, recovering
(\ref{eq:hp_tau}).

\subsection{Computational Note}

For MCMC sampling, the covariance matrix $\boldsymbol{\Sigma}$ is
parameterized via its Cholesky decomposition $\boldsymbol{\Sigma} =
\mathbf{L}\mathbf{L}^\top$, where $\mathbf{L}$ is lower triangular with
positive diagonal. The diagonal elements of $\mathbf{L}$ are stored on the log
scale to ensure positivity. This adds $p(p+1)/2$ parameters (the lower triangle
of $\mathbf{L}$) in place of the $p$ log-standard-deviation parameters of the
independent case, but guarantees that the sampled covariance matrix is always
positive definite.

% \appendixref{}  % uncomment when this appendix has citations

\pagebreak

\appsection{Prior Matching with 3 Moments: The Skew-Normal Distribution}
\label{app:skew_normal}
\titleformat*{\subsection}{\normalsize\itshape}

This appendix demonstrates the law of total cumulance for a hierarchical model
where the unconditional distribution is Skew-Normal, requiring matching of
three cumulants rather than two.

\subsection{Law of Total Cumulance for Skewness}

The law of total cumulance provides a general decomposition for the $n$-th
cumulant. For the third cumulant $\kappa_3$ (equal to the third central moment,
related to skewness):
\begin{equation}
  \label{eq:total_skewness}
  \kappa_3(Y) = \E{\kappa_3(Y|X)} + 3\Cov{\E{Y|X}}{\Var{Y|X}} + \kappa_3(\E{Y|X})
\end{equation}
\noindent The three terms have intuitive interpretations:
\begin{enumerate}
  \item $\E{\kappa_3(Y|X)}$: Expected within-group skewness
  \item $3\Cov{\E{Y|X}}{\Var{Y|X}}$: Covariance between group means and variances
  \item $\kappa_3(\E{Y|X})$: Skewness of the group means
\end{enumerate}

\subsection{The Skew-Normal Distribution}

The Skew-Normal distribution has density:
\begin{equation}
  f(x; \xi, \omega, \alpha) = \frac{2}{\omega} \phi\left(\frac{x-\xi}{\omega}\right) \Phi\left(\alpha\frac{x-\xi}{\omega}\right)
\end{equation}
where $\xi$ is location, $\omega > 0$ is scale, $\alpha$ is the shape parameter
controlling skewness, and $\phi$, $\Phi$ are the standard normal PDF and CDF.

The cumulants are:
\begin{align}
  \kappa_1 &= \xi + \omega \delta \sqrt{\frac{2}{\pi}}, \quad \text{where } \delta = \frac{\alpha}{\sqrt{1+\alpha^2}} \\
  \kappa_2 &= \omega^2 \left(1 - \frac{2\delta^2}{\pi}\right) \\
  \kappa_3 &= \frac{4-\pi}{2} \left(\delta\sqrt{\frac{2}{\pi}}\right)^3 \omega^3
\end{align}

\subsection{Hierarchical Decomposition}

Suppose we wish to specify hyperpriors such that the unconditional distribution
of $\theta_i$ matches a direct prior $\text{Skew-Normal}(\xi^*, \omega^*, \alpha^*)$.
The Skew-Normal has three parameters, so following the prior selection approach,
we must match three cumulants using the law of total cumulance.

Consider a hierarchical structure where $\theta_i \sim
\text{Skew-Normal}(\xi, \omega, \alpha)$ with all three parameters drawn from
hyperpriors. Applying the law of total cumulance by conditioning on
$(\xi, \omega, \alpha)$:

\paragraph{First cumulant (mean).}
Using the law of iterated expectations:
\begin{equation}
  \kappa_1^* = \E{\kappa_1(\theta_i | \xi, \omega, \alpha)}
             = \E{\xi + \omega \delta(\alpha) \sqrt{2/\pi}}
\end{equation}
where $\delta(\alpha) = \alpha / \sqrt{1 + \alpha^2}$.

\paragraph{Second cumulant (variance).}
Using the law of total variance:
\begin{equation}
  \kappa_2^* = \E{\kappa_2(\theta_i | \xi, \omega, \alpha)} + \Var{\kappa_1(\theta_i | \xi, \omega, \alpha)}
\end{equation}
The conditional variance is $\omega^2(1 - 2\delta(\alpha)^2/\pi)$, giving:
\begin{equation}
  \kappa_2^* = \E{\omega^2 \left(1 - \frac{2\delta(\alpha)^2}{\pi}\right)}
             + \Var{\xi + \omega \delta(\alpha) \sqrt{2/\pi}}
\end{equation}

\paragraph{Third cumulant (skewness).}
Applying (\ref{eq:total_skewness}):
\begin{equation}
  \kappa_3^* = \E{\kappa_3(\theta_i | \xi, \omega, \alpha)}
             + 3\Cov{\kappa_1(\theta_i | \xi, \omega, \alpha)}{\kappa_2(\theta_i | \xi, \omega, \alpha)}
             + \kappa_3(\kappa_1(\theta_i | \xi, \omega, \alpha))
\end{equation}
The conditional third cumulant is:
\begin{equation}
  \kappa_3(\theta_i | \xi, \omega, \alpha) = \frac{4-\pi}{2} \left(\delta(\alpha)\sqrt{\frac{2}{\pi}}\right)^3 \omega^3
\end{equation}

\paragraph{Matching conditions.}
The three cumulant equations provide constraints that link the hyperprior
parameters to the direct prior parameters $(\xi^*, \omega^*, \alpha^*)$.
Unlike the Normal case where two equations suffice for two parameters, the
Skew-Normal requires matching all three cumulants, which couples the hyperpriors
for $\xi$, $\omega$, and $\alpha$ in a more complex manner.

If independence is assumed across hyperpriors, the covariance terms vanish,
simplifying the system. However, the nonlinear dependence on $\delta(\alpha)$
means that even with independence, the matching conditions involve moments of
$\alpha$ that may not have closed-form solutions for standard hyperprior choices.

\paragraph{Hyperprior specification.}
Consider the hyperprior structure:
\begin{align}
  \xi &\sim \text{Normal}(\mu_\xi, \tau_\xi^2) \\
  \omega^2 &\sim \text{Inverse-Gamma}(\alpha_\omega, \beta_\omega) \\
  \alpha &\sim \text{Normal}(\mu_\alpha, \tau_\alpha^2)
\end{align}
This introduces 6 hyperprior parameters:
$(\mu_\xi, \tau_\xi^2, \alpha_\omega, \beta_\omega, \mu_\alpha, \tau_\alpha^2)$.
The 3 cumulant matching conditions
determine 3 of these, leaving 3 degrees of freedom for researcher choice.

\paragraph{Researcher-specified parameters.}
Following the structure from the Normal case, I propose three interpretable
parameters:
\begin{enumerate}
  \item $\rho \in (0,1)$: \textbf{Variance partition ratio.} As in the Normal
    case, this governs the fraction of prior variance attributable to
    uncertainty about population location versus population heterogeneity.

  \item $\alpha_\omega > 2$: \textbf{Scale concentration.} The shape parameter
    for the Inverse-Gamma prior on $\omega^2$, controlling confidence in the
    prior selection scale. Analogous to $\alpha_0$ in the Normal case.

  \item $\tau_\alpha^2 > 0$: \textbf{Shape uncertainty.} The prior variance on
    the population shape parameter $\alpha$. This controls how much the
    skewness can vary across the population. When $\tau_\alpha^2 \to 0$,
    skewness is fixed at the population level; larger values allow individuals
    to differ in their degree of asymmetry.
\end{enumerate}

\paragraph{Determined parameters.}
Given the researcher choices $(\rho, \alpha_\omega, \tau_\alpha^2)$ and the
direct prior cumulants $(\kappa_1^*, \kappa_2^*, \kappa_3^*)$, the remaining
three hyperprior parameters $(\mu_\xi, \beta_\omega, \mu_\alpha)$ are determined
by solving the cumulant matching conditions:
\begin{align}
  \kappa_1^* &= \mu_\xi + \sqrt{2/\pi} \cdot \E{\omega} \E{\delta(\alpha)} \\
  \kappa_2^* &= \E{\omega^2}\E{1 - 2\delta(\alpha)^2/\pi} + \tau_\xi^2
              + \frac{2}{\pi}\Var{\omega\delta(\alpha)} \\
  \kappa_3^* &= \E{\kappa_3(\theta | \xi, \omega, \alpha)}
              + \kappa_3\left(\E{\theta | \xi, \omega, \alpha}\right)
\end{align}
where $\tau_\xi^2 = \rho \cdot \kappa_2^*$ (from the variance partition), and
$\E{\omega^2} = \beta_\omega / (\alpha_\omega - 1)$.

\paragraph{Numerical solution.}
The nonlinear dependence on $\delta(\alpha) = \alpha/\sqrt{1+\alpha^2}$ generally
precludes closed-form solutions. However, given researcher choices, the system
reduces to 3 equations in 3 unknowns $(\mu_\xi, \beta_\omega, \mu_\alpha)$,
solvable via root-finding:
\begin{equation}
  \text{Find } \theta \text{such that } \kappa_n^{\text{implied}}(\theta) = \kappa_n^* \text{for } n = 1, 2, 3
\end{equation}
The implied cumulants can be computed via Monte Carlo integration over the
hyperprior distributions.

% \appendixref{}  % uncomment when this appendix has citations

\pagebreak

\appsection{MCMC Implementation Details}
\label{app:mcmc}
\titleformat*{\subsection}{\normalsize\itshape}

\subsection{Sampling Algorithm}

The model is estimated using a block-wise Metropolis-Hastings within Gibbs
sampler. The full parameter vector contains
\py{Q.get('n_total_params', 'N/A')} parameters: 4 per subject ($r$, $\ln\phi$,
$\ln\eta$, $\ln\lambda$) for \py{Q.get('n_subjects', 'N/A')} subjects, plus 8
hyperparameters (a population mean and log-standard-deviation for each of the 4
parameters).

These parameters are divided into 115 blocks:
\begin{itemize}
  \item \py{Q.get('n_subjects', 'N/A')} subject blocks, each containing the 4
    parameters for a single subject.
  \item 4 hyperparameter blocks, each containing the population mean and
    log-standard-deviation for one parameter (e.g., $\mu_r$ and $\ln\sigma_r$).
\end{itemize}

\noindent Each iteration of the sampler cycles through all 115 blocks, proposing
updates to each block conditional on the current values of all other blocks.
Subject blocks are updated via standard Metropolis-Hastings, while
hyperparameter blocks use the coupled transform sampler described below.

The sampler uses coupled A/B chain groups: half the chains are designated as
group A and the other half as group B. Cross-chain covariance estimates from
group A inform proposals for group B, and vice versa. This prevents the
proposal adaptation from introducing dependence within each group, preserving
the validity of convergence diagnostics.

\subsection{Non-Centered Parameterization}

The model uses a non-centered parameterization (NCP) to improve MCMC mixing.
In a centered parameterization, the subject-level parameters
$\theta_i$ are drawn directly from the population distribution
$\theta_i \sim N(\mu, \sigma^2)$, creating strong posterior correlations between
$\theta_i$ and the hyperparameters $(\mu, \sigma)$---the well-known
\enquote{funnel} problem.

NCP addresses this by reparameterizing in terms of standardized deviations:
\begin{align}
  \varepsilon_i &\sim N(0, 1) \\
  \theta_i &= \mu + \sigma \cdot \varepsilon_i
\end{align}

The sampler stores $\varepsilon_i$ rather than $\theta_i$. Since the
$\varepsilon_i$ are \textit{a priori} independent of $(\mu, \sigma)$, the
funnel geometry is eliminated in the transformed space, and the subject-level
blocks mix more efficiently.

\subsection{Coupled Transform Sampler}

While NCP improves the geometry for subject-level parameters, updating the
hyperparameters $(\mu, \sigma)$ via standard Metropolis-Hastings still requires
evaluating the full data likelihood at the implied
$\theta_i = \mu + \sigma \varepsilon_i$ values. When a new $(\mu', \sigma')$ is
proposed, the reconstructed $\theta_i$ values change, potentially producing a
poor fit and low acceptance rates.

The coupled transform sampler resolves this by coupling the hyperparameter
proposal with a deterministic transformation of the $\varepsilon_i$ values that
preserves the implied $\theta_i$. When proposing $(\mu', \sigma')$, the
epsilons are updated as:
\begin{equation}
  \varepsilon_i' = \frac{\theta_i - \mu'}{\sigma'}
    = \frac{\mu + \sigma \varepsilon_i - \mu'}{\sigma'}
\end{equation}

Since $\theta_i' = \mu' + \sigma' \varepsilon_i' = \theta_i$ by construction,
the data likelihood cancels from the Metropolis-Hastings acceptance ratio. The
log acceptance probability reduces to:
\begin{equation}
  \label{eq:coupled_log_alpha}
  \log \alpha =
    \underbrace{\log q(\mu, \sigma \mid \mu', \sigma')
    - \log q(\mu', \sigma' \mid \mu, \sigma)}_{\text{proposal ratio}}
    + \underbrace{N(\log\sigma - \log\sigma')}_{\text{Jacobian}}
    + \underbrace{-\tfrac{1}{2}\textstyle\sum_{i=1}^N
      (\varepsilon_i'^2 - \varepsilon_i^2)}_{\text{epsilon prior ratio}}
    + \underbrace{\log p(\mu', \sigma')
    - \log p(\mu, \sigma)}_{\text{hyperprior ratio}}
\end{equation}
where $N$ is the number of subjects and $q(\cdot \mid \cdot)$ is the proposal
distribution. The $N$-fold Jacobian factor arises because the transformation
rescales all $N$ epsilon values simultaneously. A full derivation, including the
proof of detailed balance, is provided in the supplementary technical note on
the coupled transform sampler.

This approach dramatically improves hyperparameter acceptance rates because
proposals are no longer penalized by poor likelihood fits at the new
$\theta_i$ values.

\subsection{Likelihood Function}

The likelihood for subject $i$'s choice $y_{it} \in \{A, B\}$ between lotteries
$A$ and $B$ on trial $t$ uses the Fechner (logistic) error specification from
(\ref{eq:likelihood}):
\begin{equation}
    P(y_{it} = A | \theta_i) = \frac{1}{1 + \exp\left(-\Delta_{it}\right)}
\end{equation}
where $\Delta_{it} = (U_A - U_B) / [\lambda_i (u(\bar{x}) - u(\underbar{x}))]$
as defined in (\ref{eq:likelihood}), with $\lambda_i = \exp(\log \lambda_i)$
the noise parameter.

\subsection{Computational Efficiency}

To accelerate likelihood evaluation, I precompute utilities for the 3 unique
prize amounts in the data, then use index lookups rather than recomputing these
values for each of the \py{D['n_subjects'] * D['n_choices_per_subject']}
observations.

\subsection{Proposal Adaptation}

Each block uses an adaptive Metropolis proposal that learns from cross-chain
covariance. Let $\hat{\boldsymbol{\Sigma}}_b$ denote the empirical covariance
of block $b$'s parameters across chains in the opposing A/B group. The proposal
for block $b$ is:
\begin{equation}
  \boldsymbol{\delta}_b \sim N\!\left(\mathbf{0},\;
    s_b^2 \hat{\boldsymbol{\Sigma}}_b\right)
\end{equation}
where $s_b$ is a scaling factor that is itself adapted to target an acceptance
rate near the theoretically optimal range. The cross-chain covariance is
smoothed over iterations to ensure stability, and proposals from group A use
covariance estimates from group B and vice versa to maintain the independence
needed for valid $\hat{R}$ diagnostics.

\subsection{Run Configuration}

The sampler is run with \py{Q.get('n_chains', 'N/A')} parallel chains. The
sampling proceeds through 6 sequential runs: 1 initial \enquote{reset} run
(which initializes proposal adaptation) followed by 5 \enquote{resume} runs
that continue from the previous checkpoint with accumulated adaptation
statistics. Each run collects 100 iterations with a thinning factor of 10
(i.e., every 10th sample is retained), yielding 600 total stored iterations per
chain.

The first half of the stored iterations is discarded as burn-in, leaving 300
post-burn-in iterations across \py{Q.get('n_chains', 'N/A')} chains for a
total of \py{total_samples} posterior samples used for inference.

\subsection{Convergence Diagnostics}

Convergence is assessed using the $\hat{R}$ statistic, which compares
within-chain and between-chain variance. Values near 1.0 indicate convergence;
the conventional threshold is $\hat{R} < 1.1$. Across all
\py{Q.get('n_total_params', 'N/A')} parameters, the maximum $\hat{R}$ is
\py{fmt(Q.get('rhat_max', 0), 3)}, indicating excellent convergence. Trace
plots for the hyperparameters are available in the diagnostic figures.

\subsection{Software}

Estimation uses the \texttt{bamcmc} package, a JAX-based MCMC sampler with
coupled A/B chains and GPU acceleration. JAX is configured with 64-bit
floating-point precision throughout the sampling to prevent numerical issues
in likelihood evaluation and Jacobian computations. The sampler runs
\py{Q.get('n_chains', 'N/A')} chains in parallel on an NVIDIA RTX 4070 GPU.

% \appendixref{}  % uncomment when this appendix has citations

\end{document}
